%Version 3.1 December 2024
% See section 11 of the User Manual for version history
%
%%%%%%%%%%%%%%%%%%%%%%%%%%%%%%%%%%%%%%%%%%%%%%%%%%%%%%%%%%%%%%%%%%%%%%
%%                                                                 %%
%% Please do not use \input{...} to include other tex files.       %%
%% Submit your LaTeX manuscript as one .tex document.              %%
%%                                                                 %%
%% All additional figures and files should be attached             %%
%% separately and not embedded in the \TeX\ document itself.       %%
%%                                                                 %%
%%%%%%%%%%%%%%%%%%%%%%%%%%%%%%%%%%%%%%%%%%%%%%%%%%%%%%%%%%%%%%%%%%%%%

%%\documentclass[referee,sn-basic]{sn-jnl}% referee option is meant for double line spacing

%%=======================================================%%
%% to print line numbers in the margin use lineno option %%
%%=======================================================%%

%%\documentclass[lineno,pdflatex,sn-basic]{sn-jnl}% Basic Springer Nature Reference Style/Chemistry Reference Style

%%=========================================================================================%%
%% the documentclass is set to pdflatex as default. You can delete it if not appropriate.  %%
%%=========================================================================================%%

%%\documentclass[sn-basic]{sn-jnl}% Basic Springer Nature Reference Style/Chemistry Reference Style

%%Note: the following reference styles support Namedate and Numbered referencing. By default the style follows the most common style. To switch between the options you can add or remove �Numbered� in the optional parenthesis. 
%%The option is available for: sn-basic.bst, sn-chicago.bst%  
 
%%\documentclass[pdflatex,sn-nature]{sn-jnl}% Style for submissions to Nature Portfolio journals
\documentclass[pdflatex,sn-basic]{sn-jnl}% Basic Springer Nature Reference Style/Chemistry Reference Style
%%\documentclass[pdflatex,sn-mathphys-num]{sn-jnl}% Math and Physical Sciences Numbered Reference Style
%%\documentclass[pdflatex,sn-mathphys-ay]{sn-jnl}% Math and Physical Sciences Author Year Reference Style
%%\documentclass[pdflatex,sn-aps]{sn-jnl}% American Physical Society (APS) Reference Style
%%\documentclass[pdflatex,sn-vancouver-num]{sn-jnl}% Vancouver Numbered Reference Style
%%\documentclass[pdflatex,sn-vancouver-ay]{sn-jnl}% Vancouver Author Year Reference Style
%%\documentclass[pdflatex,sn-apa]{sn-jnl}% APA Reference Style
%%\documentclass[pdflatex,sn-chicago]{sn-jnl}% Chicago-based Humanities Reference Style

%%%% Standard Packages
%%<additional latex packages if required can be included here>

\usepackage{graphicx}%
\usepackage{multirow}%
\usepackage{amsmath,amssymb,amsfonts}%
\usepackage{amsthm}%
\usepackage{mathrsfs}%
\usepackage[title]{appendix}%
\usepackage{xcolor}%
\usepackage{textcomp}%
\usepackage{manyfoot}%
\usepackage{booktabs}%
\usepackage{algorithm}%
\usepackage{algorithmicx}%
\usepackage{algpseudocode}%
\usepackage{listings}%
\usepackage{silence}%
\WarningFilter{latexfont}{Size substitutions with differences}
\WarningFilter{latexfont}{Font shape `U/rsfs/m/n' in size}
%%%%

%%%%%=============================================================================%%%%
%%%%  Remarks: This template is provided to aid authors with the preparation
%%%%  of original research articles intended for submission to journals published 
%%%%  by Springer Nature. The guidance has been prepared in partnership with 
%%%%  production teams to conform to Springer Nature technical requirements. 
%%%%  Editorial and presentation requirements differ among journal portfolios and 
%%%%  research disciplines. You may find sections in this template are irrelevant 
%%%%  to your work and are empowered to omit any such section if allowed by the 
%%%%  journal you intend to submit to. The submission guidelines and policies 
%%%%  of the journal take precedence. A detailed User Manual is available in the 
%%%%  template package for technical guidance.
%%%%%=============================================================================%%%%

%% as per the requirement new theorem styles can be included as shown below
\theoremstyle{thmstyleone}%
\newtheorem{theorem}{Theorem}%  meant for continuous numbers
%%\newtheorem{theorem}{Theorem}[section]% meant for sectionwise numbers
%% optional argument [theorem] produces theorem numbering sequence instead of independent numbers for Proposition
\newtheorem{proposition}[theorem]{Proposition}% 
%%\newtheorem{proposition}{Proposition}% to get separate numbers for theorem and proposition etc.

\theoremstyle{thmstyletwo}%
\newtheorem{example}{Example}%
\newtheorem{remark}{Remark}%

\theoremstyle{thmstylethree}%
\newtheorem{definition}{Definition}%

\newcommand{\KM}{Kaplan--Meier} %chktex 8
\newcommand{\KoMbine}{\textsc{KoMbine}}
\newcommand{\Combine}{\textsc{Combine}}
\newcommand{\RooFit}{\textsc{RooFit}}
\newcommand{\Minuit}{\textsc{Minuit}}
\newcommand{\lifelines}{lifelines}
\newcommand{\Python}{Python}
\newcommand{\scipy}{scipy}
\newcommand{\Gurobi}{Gurobi}

\raggedbottom{}
%%\unnumbered% uncomment this for unnumbered level heads

\begin{document}

\title[Propagating Statistical and Systematic Errors to \KM{} Curves]{Propagating Statistical and Systematic Errors to \KM{} Curves}

%%=============================================================%%
%% GivenName	-> \fnm{Joergen W.}
%% Particle	-> \spfx{van der} -> surname prefix
%% FamilyName	-> \sur{Ploeg}
%% Suffix	-> \sfx{IV}
%% \author*[1,2]{\fnm{Joergen W.} \spfx{van der} \sur{Ploeg} 
%%  \sfx{IV}}\email{iauthor@gmail.com}
%%=============================================================%%

\author[1,2,3]{\fnm{Jeffrey} \sur{Roskes}}\email{jroskes1@jhu.edu}

\affil[1]{\orgdiv{Institute for Data Intensive Engineering and Science}, \orgname{Johns Hopkins University}, \orgaddress{\street{3400 N. Charles Street}, \city{Baltimore}, \postcode{21218}, \state{MD}, \country{USA}}}

\affil[2]{\orgdiv{Department of Physics and Astronomy}, \orgname{Johns Hopkins University}, \orgaddress{\street{3400 N. Charles Street}, \city{Baltimore}, \postcode{21218}, \state{MD}, \country{USA}}}

\affil[3]{\orgdiv{Mark Foundation Center for Advanced Genomics and Imaging}, \orgname{Johns Hopkins University}, \orgaddress{\street{3400 N. Charles Street}, \city{Baltimore}, \postcode{21218}, \state{MD}, \country{USA}}}

%%==================================%%
%% Sample for unstructured abstract %%
%%==================================%%

\abstract{
\KM{} curves are widely used in medical research to evaluate the performance of biomarkers and predict patient outcomes. These curves are often shown without error bands, and even when error bands are provided, they typically only account for the statistical uncertainty resulting from the finite number of patients in the study. In reality, other sources of uncertainty affect the measurements as well. As datasets grow, the statistical uncertainty on the number of patients no longer dominates the overall uncertainty, and other uncertainties are increasingly important to model. We present a maximum likelihood method to propagate both statistical and systematic uncertainties through the \KM{} curve estimation processes. We illustrate our method using real-world data from patients with non-small-cell lung cancer and show that estimating and propagating these uncertainties is essential for evaluating biomarker performance.
}

%%================================%%
%% Sample for structured abstract %%
%%================================%%

% \abstract{\textbf{Purpose:} The abstract serves both as a general introduction to the topic and as a brief, non-technical summary of the main results and their implications. The abstract must not include subheadings (unless expressly permitted in the journal's Instructions to Authors), equations or citations. As a guide the abstract should not exceed 200 words. Most journals do not set a hard limit however authors are advised to check the author instructions for the journal they are submitting to.
% 
% \textbf{Methods:} The abstract serves both as a general introduction to the topic and as a brief, non-technical summary of the main results and their implications. The abstract must not include subheadings (unless expressly permitted in the journal's Instructions to Authors), equations or citations. As a guide the abstract should not exceed 200 words. Most journals do not set a hard limit however authors are advised to check the author instructions for the journal they are submitting to.
% 
% \textbf{Results:} The abstract serves both as a general introduction to the topic and as a brief, non-technical summary of the main results and their implications. The abstract must not include subheadings (unless expressly permitted in the journal's Instructions to Authors), equations or citations. As a guide the abstract should not exceed 200 words. Most journals do not set a hard limit however authors are advised to check the author instructions for the journal they are submitting to.
% 
% \textbf{Conclusion:} The abstract serves both as a general introduction to the topic and as a brief, non-technical summary of the main results and their implications. The abstract must not include subheadings (unless expressly permitted in the journal's Instructions to Authors), equations or citations. As a guide the abstract should not exceed 200 words. Most journals do not set a hard limit however authors are advised to check the author instructions for the journal they are submitting to.}

\keywords{Survival analysis, Maximum likelihood, Uncertainty, Biomarkers, Systematic uncertainty}

%%\pacs[JEL Classification]{D8, H51}

%%\pacs[MSC Classification]{35A01, 65L10, 65L12, 65L20, 65L70}

\maketitle

\section{Introduction}\label{sec:intro}

\KM{} curves are widely used in medical research to predict patient outcomes by giving the survival probability as a function of time. The curves are constructed empirically based on observed outcomes in a cohort of patients. They can be used to evaluate the performance of biomarkers, by stratifying patients into groups based on the biomarker value and plotting separate \KM{} curves for each group, or to compare the effectiveness of different treatments. By examining the \KM{} curves, we can also predict how long a future patient is likely to survive.

While it would be desireable, with high data volumes, to estimate survival probability as a function of a continuous biomarker, a more common approach is to stratify patients into discrete categories, such as ``high'' and ``low'' expression of a certain gene or density of a certain cell.  If the curves differ significantly, these biomarkers can be used to guide treatment decisions.

The biomarker measurements are subject to uncertainty, and it is desireable to incorporate this uncertainty into the \KM{} curve estimation.

Uncertainty estimates on \KM{} curves typically take into account only one source of uncertainty: the statistical uncertainty from the finite number of patients in the study.  This is typically done using \citeauthor{Greenwood} confidence intervals, which are based on the binomial distribution of deaths among at-risk patients \citep{Greenwood,GreenwoodNotes}\@.  \citeauthor{Greenwood}'s formula uses several approximations and behaves poorly in limiting cases, such as before the first death and after the last.

\Citet{FayBrittainProschan2013} survey several methods to calculate confidence intervals on \KM{} curves, including some that avoid these edge cases.  One that we wish to highlight is the method of \citet{ThomasGrunkemeier}, which uses a likelihood ratio to construct confidence intervals.  \Citeauthor{FayBrittainProschan2013} propose a beta product confidence procedure, which has better pointwise coverage for low sample sizes, and also has the desireable property of changing the error bounds when patients are censored, not only when they die.  Their method, like the others they survey, only accounts for the binomial uncertainty from the finite number of patients.

\citet{YiBook} discusses methods to handle measurement error in all kinds of regression analyses, with a chapter dedicated to survival analysis.  Most of that chapter focuses on modeling survival as a function of continuous covariates.  Section~3.7 discusses our use case: a method to handle misclassification error when using two discrete categories.

\citeauthor{YiBook}'s general strategy is to integrate over each patient's possible true group membership, using the patient's group memership probabilities, to construct a likelihood function for the hazard ratio.  Importantly, however, while the \emph{best-fit} hazard ratio estimated using her method is less biased than one estimated without accounting for measurement errors, \citeauthor{YiBook} treats the misclassification model as fixed, and, after correcting for it, considers the sampling uncertainty of the corrected estimator. Our goal is to treat the patient-wise uncertainties as well, producing confidence intervals that correctly reflect increasing uncertainty in patient assignments, growing to \([0, 1]\) in the limit of completely uninformative measurements.  The confidence intervals from \citeauthor{YiBook}'s method are not designed to reflect the patient-wise uncertainty and, as a result, do not have this property.

\citeauthor{YiBook}'s method can be naturally extended to construct weighted \KM{} curves, where each patient is weighted in the at-risk and death counts by the probability that they belong to the category being plotted.  While one could construct confidence intervals using a generalized binomial pseudo-likelihood with fractional counts, these confidence intervals would still not propagate the patient-wise uncertainties in the manner we want to do here, and would not grow with the patient-wise uncertainties, similar to the confidence intervals on the hazard ratio.

In this paper, we present, to the best of our knowledge, the first method, which we call \KoMbine{}, to propagate the uncertainties in patient group assignments to \KM{} curves.

The maximum likelihood techniques and modeling of systematic uncertainties used in our method are based on the \Combine{} package \citep{CAT-23-001}, which was developed by the CMS Collaboration in particle physics\@. \Combine, built on the \RooFit{} framework \citep{RooFit} and the \Minuit{} optimizer \citep{Minuit}, is used for high-precision measurements of particle properties, which must account for tens or even hundreds of systematic uncertainties, and was used for such high-profile measurements as the discovery of the Higgs boson \citep{HIG-12-028} and subsequent measurements of its properties.

No particle physics measurement is reported without a comprehensive assessment of the uncertainties that affect the measurement. This rigorous approach has led to a higher degree of reproducibility in particle physics results than in many other fields, including medical research.

Our method handles the statistical uncertainty on the number of patients in a more rigorous way than conventional methods and incorporates patient-wise uncertainties into the calculation. The method also allows for calculating the statistical and patient-wise uncertainties separately. Determining which uncertainties dominate can help to plan the next iteration of the experiment.

\section{Methodology}\label{sec:methodology}

\KoMbine{} uses a log-likelihood function to compute the confidence interval for each point on the \KM{} curve. The log likelihood is a function of the survival probability \(S\) as well as other parameters. We maximize the likelihood over those other parameters to find the ``likelihood scan'' as a function of \(S\) alone. The points of maximum likelihood form the \KM{} curve, and the confidence intervals are determined by the points where the log likelihood is below a threshold.

\subsection{Notation}\label{sec:notation}

We denote the inclusion of patient \(j\) by \(a_j\), which is \(1\) if the patient is included and \(0\) if the patient is excluded. The numbers of patients who were at risk and who died at time \(t_i\) are:
\begin{align}
r_i &= \sum_{t_j \geq t_i} a_j, \\
d_i &= \sum_{\substack{t_j = t_i \\ \neg c_j}} a_j,
\end{align}
where \(c_j\) indicates whether the patient was censored or not. We will allow \(a_j\) to vary so that \(r_i\) and \(d_i\) may differ from their nominal values \(\hat{r}_i\) and \(\hat{d}_i\).

The survival probability at time \(t_n\) is given by the \KM{} formula:
\begin{equation}
S_n = \prod_{i=1}^{n} p_i^\text{s}, \label{eq:km-survival}
\end{equation}
where \(p_i^\text{s}\) is the probability that a patient who has survived until time point \(t_i\) continues to live. The nominal probability is calculated from the observed survival probabilities:
\begin{equation}
\hat{S}_n = \prod_{i=1}^{n} \left(1 - \frac{\hat{d}_i}{\hat{r}_i}\right). \label{eq:km-nominal}
\end{equation}

\subsection{Greenwood confidence intervals}\label{sec:greenwood}

As mentioned earlier, confidence intervals on \KM{} curves are typically estimated using some variation of \citeauthor{Greenwood}'s formula. At confidence level \(1-\alpha\), \citeauthor{Greenwood}'s interval is given by
\begin{equation}
\label{eq:greenwood}
S(t)\in\hat{S}(t) \pm z_{\alpha/2} \sqrt{{\hat{S}(t)}^2{\sum_{i=1}^{n} \frac{d_i}{r_i(r_i-d_i)}}},
\end{equation}
where \(\hat{S}\) is the nominal value of the survival function, given by Eq.~(\ref{eq:km-nominal}). The exponential \citeauthor{Greenwood} interval is given by
\begin{equation}
\label{eq:exponential-greenwood}
\log{\left(-\log S(t)\right)} \in \log{\left(-\log \hat{S}(t)\right)} \pm z_{\alpha/2} \sqrt{\frac{1}{{\left(\log \hat{S}(t)\right)}^2} \sum_{i=1}^{n} \frac{d_i}{r_i(r_i-d_i)}}.
\end{equation}

\citet{GreenwoodNotes} gives the derivation of these formulas, which makes certain approximations: the ``delta-method'' approximation, which requires that the true and observed probabilities are close enough that \(\log\left(-\log S(t)\right)\) can be treated as approximately linear, and the approximation of the binomial distribution by a normal distribution.

With large numbers of patients, these approximations hold, but their validity decreases for smaller cohort sizes.  Several of the methods surveyed by \citet{FayBrittainProschan2013} avoid these approximations and behave better with a small number of patients and in limiting cases, such as when all patients survive or all patients die, as shown in Fig.~\ref{fig:compare-to-greenwood}.

\begin{figure}
\centering
\includegraphics[width=\linewidth]{comparison_to_greenwood.pdf}
\caption{\label{fig:compare-to-greenwood} \KM{} curves with confidence intervals using only the binomial error, calculated using the exponential \citeauthor{Greenwood} (e. G.) formula from Eq.~(\ref{eq:exponential-greenwood}) (gold) and \citeauthor{ThomasGrunkemeier}'s method implemented in \KoMbine{} (blue). Panel A shows a small cohort of 12 patients, while panel B shows a larger cohort of 100 patients.}
\end{figure}

We choose to extend the likelihood method of \citet{ThomasGrunkemeier}.  Although \citeauthor{FayBrittainProschan2013}'s method has some advantages over \citeauthor{ThomasGrunkemeier}'s, our goal is to incorporate patient-wise uncertainty, and the likelihood method provides a natural way to do so.

\subsection{Log likelihood method}

For clarity, we will distinguish the two types of uncertainty: the statistical uncertainty from the finite number of patients (which we will call the ``binomial uncertainty'') and the patient-wise uncertainty.

\subsubsection{Binomial uncertainty}\label{sec:binomial-uncertainty}

The binomial contribution to the \KM{} curve is computed from the number of patients at risk and the number of patients who died at each time point (\citealp{ThomasGrunkemeier}; see also \citealp{Li1995}, who shows that this represents a genuine nonparametric likelihood ratio):
\begin{equation}
\mathcal{L}_{\text{binomial}}(p_i^\text{s}, a_j) = \prod_{i=1}^{n} \binom{r_i}{d_i} {\left(p_i^\text{d}\right)}^{d_i} {\left(p_i^\text{s}\right)}^{r_i-d_i},
\end{equation}
where \(p_i^\text{d} = 1 - p_i^\text{s}\). The negative log likelihood (NLL) is
\begin{equation}
-\log \mathcal{L}_{\text{binomial}}(p_i^\text{s}, a_j) = -\left[\sum_{i=1}^{n} \left( \log\binom{r_i}{d_i} + d_i \log p_i^\text{d} + (r_i-d_i) \log p_i^\text{s} \right)\right]. \label{eq:binomial-nll}
\end{equation}

We are not interested in the individual \(p_i^\text{s}\) values, but rather the overall survival probability \(S_n\). We therefore use Eq.~(\ref{eq:km-survival}) as a constraint when we minimize the NLL\@.

\subsubsection{Patient-wise uncertainty}\label{sec:patient-wise-uncertainty}

We now extend \citeauthor{ThomasGrunkemeier}'s likelihood and consider that the patients themselves may be subject to uncertainty. Some patients who were nominally included in the \KM{} curve may actually be excluded or vice versa.

The classic case when this happens is when dividing the patients into groups based on a continuous parameter, such as gene expression or the density of a certain cell type. The parameter measurement is subject to uncertainty, which must be considered when determining the overall uncertainty on the \KM{} curve. (This is in contrast to \KM{} curves based on discrete categories with no error, such as treatment vs.\ placebo, where patient-wise uncertainty likely does not enter.)

We first precompute the NLL penalty for each patient to be included in the \KM{} curve, which is negative if the patient is nominally included and positive if the patient is nominally excluded. The penalty depends on the functional form of the patient's parameter's probability density function, and it is calculated as the difference between the minimum NLL for the parameter to be within the range to be included in the curve and the minimum NLL for it to be outside the range.

The current \KoMbine{} implementation allows several possible probability distributions for the patient parameters, including a delta function, Poisson distribution, a Poisson distribution divided by a fixed area (to represent, for instance, a density of cells), and a ratio of two Poisson distributions. Any of these can be multiplied by systematic uncertainties following the log-normal distribution, a common choice used to represent multiplicative uncertainties.

For the purpose of the math here, all that is important is that the penalty is a fixed number per patient, which we call \(-\log \mathcal{L}_j^{\text{patient}}\). We compute the NLL as
\begin{equation}
-\log \mathcal{L}_{\text{patient}}(a_j) = -\sum_{j=1}^{m} a_j\log \mathcal{L}_j^{\text{patient}}. \label{eq:patient-nll}
\end{equation}

\subsubsection{Combining the uncertainties}\label{sec:combining-uncertainties}

The total NLL is the sum of the binomial NLL and the patient-wise NLL\@:
\begin{equation}
-\log \mathcal{L}_{\text{total}}(p_i^\text{s}, a_j) = -\log \mathcal{L}_{\text{binomial}}(p_i^\text{s}, a_j) - \log \mathcal{L}_{\text{patient}}(a_j). \label{eq:total-nll}
\end{equation}

We minimize the overall NLL (\ref{eq:total-nll}) over the \(p_i^\text{s}\) and \(a_j\), with the constraint (\ref{eq:km-survival}). This is a mixed integer nonlinear programming (MINLP) problem.

\subsection{Isolating the binomial and patient-wise uncertainties}\label{sec:isolating-uncertainties}

Although the total uncertainty on the \KM{} curve includes both the binomial and the patient-wise uncertainty, it is instructive to look at them separately. Determining which uncertainty dominates can help to understand the next steps in the analysis: should we try to enroll more patients or should we improve the biomarker measurements?

\subsubsection{Binomial uncertainty}

To examine the binomial uncertainty, we can fix the parameters to their nominal values, fixing \(a_j=\hat{a}_j=1\) for the patients who are nominally included and \(0\) for those who are nominally excluded. This is equivalent to \citeauthor{ThomasGrunkemeier}'s method.

\subsubsection{Patient-wise uncertainty}

Isolating the patient-wise uncertainty is somewhat nontrivial. We cannot allow the survival probabilities \(p_i^\text{s}\) to vary freely, as we would need the binomial term to constrain them. Instead, we fix
\begin{equation}
p_i^\text{s}(a_j) = 1 - \frac{d_i(a_j)}{r_i(a_j)}.
\end{equation}

With this definition, \(S\) can only take on certain fixed, discrete values. Furthermore, depending on the order that patients were censored or died, values of \(S\) that are close together may actually only be reachable by very different combinations of patients included in the analysis, with correspondingly different NLL values. A scatter plot of \(-\log\mathcal{L}\) as a function of \(S\) would be useless.

Instead, we use a continuous variable \(S'\). We calculate \(\hat{S}\) using the nominal values of the patient parameters, as in Eq.~(\ref{eq:km-nominal}). Then, for any value of \(S'\), we minimize the NLL over \(a_j\) with the constraint that \(S\) is at least as far from \(\hat{S}\) as \(S'\):
\begin{align}
S \ge S' & \text{ if } S' > \hat{S} \\
S \le S' & \text{ if } S' < \hat{S}
\end{align}
This gives an interpretable visualization of the patient-wise uncertainty: we obtain the NLL that the survival probability is at least as far from the nominal probability as \(S'\). By construction, \(-\log\mathcal{L}(S')\) has a single minimum (either a single point or, if some patients have a penalty of \(0\), a flat region) at \(S'=\hat{S}\) and increases as \(S'\) moves away in either direction.

\subsection{Hazard ratio}

We can also use a similar method to compute the hazard ratio and its confidence interval.  We have two curves, labeled \(k=0\) and \(k=1\), and we add this index to all relevant quantities. For example, each patient may be in one of the curves, and this inclusion is given by \(a_{jk}\). Each patient can be in at most one of the two curves, but we also allow for a patient to be in neither. (This is useful, for example, when stratifying patients into low, middle, and high categories based on a parameter\@. If we compare the curves pairwise, then in each comparison, some patients will be in neither curve.)  Each patient has a positive or negative penalty, \(-\log\mathcal{L}_{jk}^\text{patient}\), to be included in each curve.

Each curve has its probabilities of dying \(p^\text{d}_{ik}\). We take the proportional hazard model,
\begin{equation}
p^\text{d}_{i0} = H p^\text{d}_{i1},
\label{eq:proportional-hazards}
\end{equation}
for a constant hazard ratio \(H\) across all time points.

The NLL for the hazard ratio is similar to Eq.~(\ref{eq:total-nll}). We have two separate patient-wise NLL penalties, one for each curve. The probability distribution representing the relationship between the curves is modeled with the Cox partial likelihood distribution with \(\beta=\log{H}\), which reduces to the hypergeometric distribution when \(H=1\)~\citep{cox1972regression,breslow1974covariance}:
\begin{align}
-\log \mathcal{L}_{\text{Cox}}(a_{jk}, H)&=-\sum_i\left[\log\binom{r_{i0}}{d_{i0}}+\log\binom{r_{i1}}{d_{i1}}+d_{i0}\log H-\log Z_i(H)\right] \label{eq:cox-nll-exact}\\
&\approx-\sum_i \left[d_{i0} \log H - d_i \log\left(H r_{i0} + r_{i1}\right)\right] \label{eq:cox-nll-breslow}
\end{align}

Eq.~(\ref{eq:cox-nll-exact}) gives the exact formula, with the normalization factor \(Z_i(H)\) given by
\begin{equation}
Z_i(H) = \sum_{m}\binom{r_{i0}}{m}\binom{r_{i1}}{d_{i0}+d_{i1}-m}H^m,
\label{eq:cox-normalization}
\end{equation}
where the sum is over all feasible values of \(m\) that allow both binomial coefficients to be finite.

Eq.~(\ref{eq:cox-nll-breslow}) applies the Breslow approximation~\citep{breslow1974covariance}. Both give identical results when no two patients die at the same time. Since our biomedical analyses typically do not involve many patients, ties are rare, and we therefore use Eq.~(\ref{eq:cox-nll-breslow}) in \KoMbine{}. We explore the effect of this approximation in Section~\ref{sec:compare-to-conventional-p-value}.

The total NLL is then given by
\begin{equation}
-\log \mathcal{L}_{\text{total}}(a_{jk}, H) = -\log\mathcal{L}_\text{Cox}(a_{jk}, H) - \sum_{k}\log\mathcal{L}_k^{\text{patient}}(a_{jk}).
\label{eq:hazard-ratio-nll}
\end{equation}
We find the hazard ratio \(H\) by minimizing the NLL over \(a_{jk}\) and \(H\). We then compute its confidence interval by scanning the NLL as a function of \(H\) and finding the points where it crosses thresholds derived from the \(\chi^2\) distribution with one degree of freedom.

\subsection[p~value]{\(p\)~value}

Building on the hazard ratio calculation, we can compute the \(p\)~value to reject the null hypothesis that two survival curves are identical. In \KoMbine{}, we use a log-rank approach, similar to the conventional approach used in \lifelines{} and similar packages. However, as in the case of the \KM{} error bars, our likelihood modeling does not rely on approximating the binomial distribution as normal, and we also include the patient-wise errors.

The conventional approach defines the residual for each time point as
\begin{align}
\begin{aligned}
U_i&=d_{i1} - d_{i1}^\text{exp} \\
&=d_{i1} - (d_{i0} + d_{i1})\frac{n_{i1}}{n_{i0}+n_{i1}}.
\end{aligned}
\end{align}
By symmetry, if we swapped the two curves, \(U_i\rightarrow-U_i\). Summing the residuals and dividing by the variance gives us the test statistic,
\begin{align}
U&=\sum_i U_i \\
Q&=\frac{U^2}{\text{Var}(U)},
\end{align}
which follows a \(\chi^2\) distribution with one degree of freedom~\citep[pp.~45--47]{mantel1966evaluation,peto1972asymptotically,appliedsurvivalanalysis}.

In the likelihood approach, we use Eq.~(\ref{eq:hazard-ratio-nll}). The null hypothesis is that \(H=1\), meaning that the two curves are identical. The alternative hypothesis allows \(H\) to float.

We evaluate the NLL for the null hypothesis by minimizing Eq.~(\ref{eq:hazard-ratio-nll}) over \(a_{jk}\) with the constraint \(H=1\). For the alternate hypothesis, we minimize Eq.~(\ref{eq:hazard-ratio-nll}) over both \(a_{jk}\) and \(H\). The difference between the two NLLs is used to compute the \(p\)~value using a \(\chi^2\) distribution with one degree of freedom.

We also allow fixing \(a_{jk}\) and using only the Cox penalty. Because floating \(a_{jk}\) can decrease the NLL for both the null and alternate hypotheses, the \(p\)~value may be either larger or smaller than the \(p\)~value with fixed \(a_{jk}\). Section~\ref{sec:compare-to-conventional-p-value} compares our Cox-only \(p\)~value to the conventional log-rank test.

\section{Numerical examples}

These examples illustrate interesting aspects of the \KM{} curve estimation in the \KoMbine{} method.

\subsection%
[Best fit vs. nominal \KM{} curve]% chktex 12
{Best fit vs.\ nominal \KM{} curve}

\begin{figure}
  \centering
  \includegraphics[width=0.5\textwidth]{km_example.pdf}
  \caption{\label{fig:best-fit-vs-nominal} An example \KM{} curve. The best fit (blue) does not coincide with the nominal curve (red). The error bands are shown as shaded blue areas.}
\end{figure}

Counter-intuitively, the best fit \KM{} curve does not necessarily coincide with the nominal \KM{} curve. An example is shown in Fig.~\ref{fig:best-fit-vs-nominal}.

If only the binomial penalty is included, the survival NLL always minimizes at the nominal survival probability. Similarly, with only the patient-wise errors, the nominal survival probability is also always optimal. However, when both are included, the two may differ.

To understand this, imagine a case where 7 patients were at risk at a particular time point and 4 of them survived. An additional patient who survived is nominally excluded from the curve, but their parameter is so close to the boundary that they can be included with only a small NLL penalty of \(0.01\)\@. Equation.~(\ref{eq:binomial-nll}) evaluates to \(-1.255\) for \(4\) patients surviving out of \(7\) with \(p_\text{s}=4/7\), but to \(-1.267\) for \(5\) patients surviving out of \(8\) with \(p_\text{s}=5/8\). This small difference in the binomial penalty outweighs the even smaller penalty for including the patient, and so the minimum NLL occurs at the best fit survival probability of \(5/8\).

This surprising result highlights the importance of the patient-wise uncertainties. If the NLL penalty for including the patient really is that small, then we cannot, with any confidence, assign the patient to one \KM{} curve or the other.

\subsection[Conventional log-rank p~value estimation]{Conventional log-rank \(p\)~value estimation}\label{sec:compare-to-conventional-p-value}

In order to compare the \(p\)~values, we generate Monte Carlo cohorts of patients. We calculate the \(p\)~values using our method, disabling the patient-wise penalties, and with the conventional log-rank test. We plot the results in Fig.~\ref{fig:compare-p-value}. The top row allows patients to die at the same time, which means that we rely on the Breslow approximation in Eq.~(\ref{eq:cox-nll-breslow}). In the bottom row, no two patients die at the same time, so the Breslow approximation is identical to the exact formula in Eq.~(\ref{eq:cox-nll-exact}).

\begin{figure}[ht]
\centering
\includegraphics[width=\linewidth]{p_value_comparison.pdf}
\caption{\label{fig:compare-p-value} Scatter plot showing the \(p\)~values computed using the conventional log-rank method (horizontal axis) and the likelihood method using only the Cox penalty (vertical axis). Each point represents a trial with randomly simulated patients. Panels A, B, and C show cohorts of 10, 100, and 1000 patients, respectively, allowing patients to die at the same time (relying on the Breslow approximation). Panels D and E show cohorts of 10 and 100 patients, respectively, with no two patients dying at the same time (so the Breslow approximation is identical to the exact formula).}
\end{figure}

When there are no ties, the points scatter along the diagonal line \(y=x\), meaning that, on average, the two methods give identical results. When ties are allowed, the diagonal line is slightly altered to a curve. The Breslow approximation gives more conservative \(p\)~values, and the effect is most pronounced for small \(p\)~values.

The Monte Carlo patients for the top row are simulated with integer survival times from \(1\) to \(10\), which means that we expect a tie rate of about \(10\%\). The error introduced by the Breslow approximation will increase with a higher tie rate.

The scatter from the line or curve, which shrinks with larger cohort sizes, is because the full likelihood method avoids the asymptotic approximations used in the conventional method, similar to the difference between the \citeauthor{ThomasGrunkemeier} and \citeauthor{Greenwood} confidence intervals (Fig.~\ref{fig:compare-to-greenwood}).

\section{Application to Non-Small Cell Lung Cancer}

\begin{figure}[ht]
  \centering
  \includegraphics[width=\linewidth]{lung_km_RFS.pdf}
  \caption{\label{fig:lung-dataset} \KM{} curves for the regression-free survival (RFS) of patients with non-small-cell lung cancer who were treated with neoadjuvant anti-PD-1 immunotherapy, stratifying patients by their pre-treatment densities of either CD8+FoxP3+ cells (left side, panels A--E) or DONUTS (right side, panels F--J).  The top row shows statistical uncertainties from the finite number of patients (panels A and F) and from the number of cells or DONUTS (panels B and G).  The second row shows the systematic uncertainty from flatfielding, before (panels C and H) and after (panels D and I) applying the corrections.  The larger plots at the bottom (panels E and J) show the combined uncertainty: both statistical uncertainties and the remaining flatfielding uncertainty after the correction was applied.} %chktex 8
\end{figure}

This example is taken from \citet{DONUTS}, in which we analyzed pre-treatment biopsies from a cohort of patients with non-small-cell lung cancer (NSCLC) treated with neoadjuvant anti-PD-1 immunotherapy, scanned on an Akoya Vectra~3.0 Imaging System. We identified CD8+FoxP3+ cells, which are associated with good outcomes, but are extremely rare: a patient may have only one or two in their biopsy. We also developed a probabilistic biomarker, called DONUTS, which resemble the niche of the CD8+FoxP3+ cells but are much more abundant.

In further work \citep{SITC2025}, we examined one systematic uncertainty that affects our measurement.  The microscope scans the tissue by dividing it into tiles, or high-power fields (HPFs), and imaging each one separately.  Due to imperfections in the imaging system, the illumination varies across the microscope's field of view.  This effect can be measured by averaging the illumination over thousands of HPFs, and we correct for it in the AstroPath pipeline \citep{astropathscience}.

This ``flatfield'' effect has a measurable impact on the cell and DONUTS densities.  In order to examine the importance of applying this correction as well as estimate the impact of any residual error, we processed the data through the AstroPath pipeline again, skipping the flatfield correction step.  We found that, while the overall densities of CD8+FoxP3+ cells and DONUTS changed only slightly, 20\% of CD8+FoxP3+ cells and 13\% of DONUTS from the original analysis were no longer identified when not applying the flatfielding correction, while 20\% new CD8+FoxP3+ cells and 13\% new DONUTS appeared.  This indicates that we can estimate that we would have a 20\% uncertainty on the cells and a 13\% uncertainty on the DONUTS if we did not apply the flatfielding correction.

For the Vectra~3.0 microscope, the pixel standard deviation across an HPF is 3.45\% before the correction and 0.61\% afterwards: approximately an 80\% reduction.  To first order, we can therefore estimate that, even after applying the flatfielding correction, 20\% of the original error is still present: a 4\% uncertainty on the CD8+FoxP3+ cells and a 2.6\% uncertainty on the DONUTS\@.  This error comes from residual imperfections in the flatfielding correction.

We used KoMbine to propagate three categories of errors to \KM{} curves for regression-free survival (RFS)\@: the binomial uncertainty from the finite number of patients, the Poisson uncertainty from the finite number of cells or DONUTS, and the systematic uncertainty from the flatfielding, calculated both before and after applying the correction.  \KM{} curves with each of these uncertainty components are shown in Fig.~\ref{fig:lung-dataset}.  The bottom plots show the combined uncertainties.

The binomial uncertainties (A and F) are similar for the cells and for the DONUTS\@, as both analyses have similar numbers of patients in the categories. The Poisson uncertainties (B and G), on the other hand, have a similar magnitude to the binomial uncertainties for the CD8+FoxP3+ cells, but are much smaller for the DONUTS, which are much more abundant.

The systematic uncertainties from flatfielding, whether before or after correction, are small for the CD8+FoxP3+ cells, even before correction.  This is because even a 20\% uncertainty on the very small number of cells does not significantly affect the patient stratification: in fact, patients with \emph{any} CD8+FoxP3+ cells in their biopsy are all assigned to the high-density group.  For the DONUTS, before correction, the flatfielding uncertainty is comparable to the binomial uncertainty.  After correction, it is subdominant and similar to the Poisson uncertainty.

It should be noted that this is a fairly small cohort of only 25 patients. In a larger-scale study, the binomial uncertainties would be smaller, while the difference in the patient-wise uncertainties would remain the same.  The difference in the overall uncertainty would therefore be even more pronounced.

\section{Discussion}

In \KoMbine{}, we introduce a novel likelihood-based method for the propagation of statistical and systematic uncertainties through the \KM{} survival function. Our handling of the binomial uncertainty is more rigorous than the conventional \citeauthor{Greenwood} confidence intervals, and the method can handle patient-wise uncertainties for the first time.

In addition to providing more precise uncertainty estimates, our method can provide critical insight for experimental design.  Researchers can avoid wasting time and effort trying to reduce subdominant uncertainties and instead focus on their work on the parts of the experiment that have the highest impact on the results.

A primary limitation of our current framework is that it is limited to uncorrelated systematic errors.  Extending the methodology to uncertainties that are correlated between patients, such as batch effects, would involve modifying our minimization problem structure to float the uncertainty parameters and computing \(a_j\) from them, rather than precomputing penalties and floating \(a_j\) directly.  Although straightforward to describe, this conceptual switch is nontrivial to implement within the current optimization framework described in Appendix~\ref{sec:implementation-details} and will probably require a rethinking of the optimization strategy.

\section{Conclusions}

Data sets used in biomedical analyses are growing rapidly, reducing statistical uncertainties.  The systematic uncertainties in the analysis are increasingly the dominant component of the overall uncertainty, and must be rigorously incorporated in order to ensure the validity of the results.

\KoMbine{} is the first methodology to formally propagate both statistical and systematic uncertainties to \KM{} survival curves.  By calculating the uncertainty budget as well as the overall confidence bands, \KoMbine{} will facilitate the development of the next generation of robust, reproducible, and clinically actionable biomarkers.

\backmatter%
\section*{Acknowledgements}

I would like to thank my AstroPath colleagues, in particular my mentors, Drs.~Alex Szalay and Janis Taube, for their support on this project. I would also like to acknowledge editing support by Ben Cohen, provided by the Office of the Vice Provost for Research at Johns Hopkins University.

\section*{Declarations}

\subsection*{Funding}
This research was supported by the Mark Foundation for Cancer Research.

\subsection*{Conflict of interest}
The author holds a patent on the DONUTS biomarker and has received licensing royalties as well as research support from Quanterix.

\subsection*{Ethics approval}
Not applicable.

\subsection*{Consent to participate}
Not applicable.

\subsection*{Consent for publication}
Not applicable.

\subsection*{Data availability}
The data used in this work are available together with the code at \url{https://github.com/AstroPathJHU/KoMbine}.

\subsection*{Code availability}
All code used in this work is available at \url{https://github.com/AstroPathJHU/KoMbine}, or at PyPI via \texttt{pip install kombinekm}.

\subsection*{Authors' contributions}
Not applicable. The author is solely responsible for all aspects of this work.

\begin{appendices}

\section{Implementation details}\label{sec:implementation-details}

\KoMbine{} is implemented in \Python{} and uses \Gurobi{} \citep{gurobi} to minimize the NLL for a given value of \(S\).

We first minimize \(-\log \mathcal{L}_{\text{total}}\) without the constraint on \(S\) (\ref{eq:km-survival}) to find the maximum likelihood estimate and the best value of \(S\)\@. \texttt{scipy.optimize.brentq} \citep{brentq,scipy} is then used to compute the confidence intervals, bounded by the values of \(S\) where the NLL equals a threshold value (1.0 for 68\% CL and 3.84 for 95\% CL).

For the patient-wise-only error, there is no need to minimize the NLL over \(S'\), as by construction it minimizes at \(S'=S_\text{nominal}\). To find the confidence intervals, we use a simple binary search that only looks at the possible discrete values of \(S\): intermediate values of \(S'\) cannot return anything different.

To compute the \(p\)~value, \KoMbine{} similarly uses \Gurobi{} to minimize the NLL for the null and alternate hypotheses. The \(p\)~value is computed from the difference between the NLLs, which follows a \(\chi^2\) distribution, using standard functions in \scipy{} \citep{scipy}.

\subsection{Collapsing consecutive deaths}

The binomial term in the NLL [Eq.~(\ref{eq:binomial-nll})] requires a sum over all time points where at least one patient died, which can be slow when there are many such time points. A mathematical trick allows us to collapse multiple time points into one, as long as no patients are censored in between, without changing the results.

Take the case of two consecutive time points, \(t_1\) and \(t_{2}\), with death counts \(d_1\) and \(d_{2}\) and risk sets \(r_1\) and \(r_{2}=r_1-d_1\)\@.  The contribution to the NLL from these two time points is
\begin{align}
\begin{aligned}
-\log \mathcal{L}_{\text{binomial},1,2} =&-\left[\log\binom{r_1}{d_1} + d_1 \log \left(1-p_1^\text{s}\right) + (r_1-d_1) \log p_1^\text{s}\right] \\
&-\left[\log\binom{r_{1}-d_{1}}{d_{2}} + d_{2} \log \left(1-p_{2}^\text{s}\right) + (r_{1}-d_{1}-d_{2}) \log p_{2}^\text{s}\right].
\end{aligned}
\label{eq:two-timepoints-nll}
\end{align}

The key observation is that \(p_1^\text{s}\) and \(p_2^\text{s}\) only enter the NLL in this term and in the constraint Eq.~(\ref{eq:km-survival}), which only depends on their product \(p_{1,2}^\text{s}=p_1^\text{s} p_2^\text{s}\)\@. We can therefore minimize Eq.~(\ref{eq:two-timepoints-nll}) over \(p_1^\text{s}\) and \(p_2^\text{s}\) with the constraint \(p_{1,2}^\text{s}=p_1^\text{s} p_2^\text{s}\)\@.  The result is
\begin{align}
p_1^\text{s} &= \frac{d_1 p_{1,2}^\text{s}+d_2}{d_1+d_2}, \\
p_2^\text{s} &= \frac{(d_1+d_2) p_{1,2}^\text{s}}{d_1 p_{1,2}^\text{s} + d_2}.
\end{align}

We plug into Eq.~(\ref{eq:two-timepoints-nll}) to obtain:
{
\newcommand{\firstline}{\log\binom{r_1}{d_1+d_2} + (d_1+d_2) \log\left(1-p_{1,2}^\text{s}\right)}
\newcommand{\secondline}{+ (r_1-d_1-d_2) \log{}p_{1,2}^\text{s}}
\begin{align}
\begin{aligned}
-\log \mathcal{L}_{\text{binomial},1,2} =&-\left[\firstline\vphantom{\secondline}\right. \\
&\qquad\qquad\qquad\qquad\qquad\qquad\qquad\hspace{0.013\textwidth}\left.\vphantom{\firstline}\secondline\right] \\
&-\left[\log\binom{d_1+d_2}{d_1} + d_1 \log \left(\frac{d_1}{d_1+d_2}\right) + d_2 \log \left(\frac{d_2}{d_1+d_2}\right)\right],
\end{aligned}
\label{eq:two-timepoints-collapsed-nll}
\end{align}
}
making use of the binomial coefficient identity
\begin{equation}
\binom{r_1}{d_1}\binom{r_1-d_1}{d_2} = \binom{r_1}{d_1+d_2}\binom{d_1+d_2}{d_1}.
\end{equation}

The first line of Eq.~(\ref{eq:two-timepoints-collapsed-nll}) is identical to the NLL contribution from a single time point with \(d_1+d_2\) deaths and risk set \(r_1\)\@. The second line does not depend on \(p_{1,2}^\text{s}\).

This procedure can be applied recursively to collapse any number of consecutive time points with no censoring in between. The result is a single binomial term for the collapsed time points, plus a multinomial NLL term for the number of patients who died at each of the original time points out of the total number of patients who died in the group:
\begin{align}
\begin{aligned}
-\log \mathcal{L}_{\text{binomial},1 \cdots n} =&-\left[\log\binom{r_1}{d_{1 \cdots n}} + d_{1 \cdots n} \log \left(1-p_{1 \cdots n}^\text{s}\right) + (r_1-d_{1 \cdots n}) \log p_{1 \cdots n}^\text{s}\right] \\
&-\left[\log{(d_{1 \cdots n}!)}-d_{1 \cdots n}\log d_{1 \cdots n} - \left(\sum_{i=1}^n \left[\log{(d_i!)} - d_i \log d_i\right]\right)\right],
\end{aligned}
\end{align}
with \(d_{1 \cdots n} = d_1 + d_2 + \cdots + d_n\) and \(p_{1 \cdots n}^\text{s} = p_1^\text{s} p_2^\text{s} \cdots p_n^\text{s}\)\@.

\KoMbine{} automatically collapses such time points before performing the NLL minimization. When applying binomial-only errors, the second line only contributes a constant offset to the NLL\@. Although it does affect the NLL when patient-wise errors are included, it can be precomputed for the various possible values of \(d_i\) and implemented in \Gurobi{} using indicator variables, thereby reducing the number of continuous variables. This can significantly speed up the NLL minimization, especially when not too many patients are censored.

\subsection{Generative AI usage}

Some of the code within the KoMbine package, particularly the testing and plotting code and the command line interface, was generated using ChatGPT, Gemini, and GitHub Copilot. All AI-generated code was thoroughly reviewed and verified by the author to ensure correctness, reproducibility, and compliance with the statistical methodology presented in Section 2.

No text of this manuscript was generated by AI, beyond AI-assisted copy editing.

\end{appendices}

%%===========================================================================================%%
%% If you are submitting to one of the Nature Portfolio journals, using the eJP submission   %%
%% system, please include the references within the manuscript file itself. You may do this  %%
%% by copying the reference list from your .bbl file, paste it into the main manuscript .tex %%
%% file, and delete the associated \verb+\bibliography+ commands.                            %%
%%===========================================================================================%%

\bibliography{02_kombine}% common bib file
%% if required, the content of .bbl file can be included here once bbl is generated
%%\input sn-article.bbl

\end{document}
